% ========================================================
%  《习近平新时代中国特色社会主义思想概论》全真模拟试卷（含参考答案）
%  覆盖范围：导论、第一章至第十二章核心考点
% ========================================================
\documentclass[12pt, a4paper]{article}

% ---- 中文支持 ----
\usepackage[UTF8]{ctex}

% ---- 数学 ----
\usepackage{amsmath, amssymb, amsthm}
\usepackage{mathtools}

% ---- 页面布局 ----
\usepackage[top=2.5cm, bottom=2.5cm, left=2.8cm, right=2.8cm]{geometry}

% ---- 表格 ----
\usepackage{booktabs, array, multirow, makecell}
\usepackage{tabularx}

% ---- 页眉页脚 ----
\usepackage{fancyhdr}
\usepackage{lastpage}
\pagestyle{fancy}
\fancyhf{}
\renewcommand{\headrulewidth}{0.6pt}
\lhead{\small 习近平新时代中国特色社会主义思想概论·GLM·A卷}
\rhead{\small 共 \pageref{LastPage} 页}
\cfoot{\small 第 \thepage\ 页}

% ---- 计数器与样式 ----
\usepackage{enumitem}
\usepackage{xcolor}
\usepackage{tcolorbox}
\tcbuselibrary{skins, breakable}

% ---- 填空线 ----
\newcommand{\blank}[1]{\underline{\hspace{#1}}}

% ---- 答案盒子 ----
\newtcolorbox{answerbox}[1][]{
  breakable,
  colback=gray!6,
  colframe=gray!50,
  fonttitle=\bfseries\small,
  title=参考答案,
  #1
}

% ---- 题号格式 ----
\newcounter{bigq}
\newcommand{\BigQ}[1]{%
  \stepcounter{bigq}%
  \vspace{6pt}%
  \noindent{\large\bfseries 第\chinese{bigq}大题\quad #1}%
  \vspace{4pt}\par
}

\newcounter{subq}[bigq]
\newcommand{\SubQ}{%
  \stepcounter{subq}%
  \noindent\textbf{（\arabic{subq}）}\quad
}

% 分隔线
\newcommand{\examrule}{\noindent\rule{\linewidth}{0.6pt}}

% 单选题选项格式（缩进+紧凑排列）
\newcommand{\opts}[4]{% 
  \par\noindent\vspace{2pt}%
  \hspace*{2em}\begin{minipage}[t]{\dimexpr\linewidth-2em\relax}
    \setlength{\parskip}{2pt}
    A.\ #1\par
    B.\ #2\par
    C.\ #3\par
    D.\ #4
  \end{minipage}\par\vspace{4pt}
}

% 多选题选项格式
\newcommand{\mopts}[4]{%
  \par\noindent\vspace{2pt}%
  \hspace*{2em}\begin{minipage}[t]{\dimexpr\linewidth-2em\relax}
    \setlength{\parskip}{2pt}
    A.\ #1\par
    B.\ #2\par
    C.\ #3\par
    D.\ #4
  \end{minipage}\par\vspace{4pt}
}

% ---- 文档开始 ----
\begin{document}

% ============================================================
%  封面
% ============================================================
\begin{center}
  {\LARGE\bfseries 习近平新时代中国特色社会主义思想概论}\\[8pt]
  {\large\bfseries 全真模拟仿真试卷（含参考答案）}\\[16pt]
  \begin{tabular}{lll}
    考试时间：120 分钟 & \quad & 满分：100 分 \\[4pt]
    \multicolumn{3}{l}{注意事项：请将答案写在答题纸指定区域，写在本卷上无效。}
  \end{tabular}
  \vspace{4pt}

  \examrule

  \vspace{6pt}
  \begin{tabular}{|c|c|c|c|c|c|}
    \hline
    \textbf{题号} & 一 & 二 & 三 & \textbf{总分} & \textbf{复核人} \\
    \hline
    \textbf{满分} & 30 & 20 & 50 & 100 & \\
    \hline
    \textbf{得分} & & & & & \\
    \hline
  \end{tabular}
\end{center}

\vspace{10pt}
\examrule
\vspace{6pt}

% ============================================================
%  第一大题：单项选择题（30 题，每题 1 分，共 30 分）
% ============================================================
\BigQ{单项选择题（每题 1 分，共 30 分。每题只有一个正确选项）}

\SubQ 习近平新时代中国特色社会主义思想创立的时代背景之一是（\quad）。
\opts{世界百年未有之大变局加速演进}{经济全球化出现逆转}{冷战思维重新主导国际关系}{西方现代化模式主导世界进程}

\SubQ "两个确立"是指：确立习近平同志党中央的核心、全党的核心地位，确立（\quad）的指导地位。
\opts{邓小平理论}{"三个代表"重要思想}{科学发展观}{习近平新时代中国特色社会主义思想}

\SubQ 中国特色社会主义进入新时代，是我国发展新的（\quad）。
\opts{历史方位}{历史阶段}{历史时期}{历史节点}

\SubQ 中国最大的国情是（\quad）。
\opts{人口基数大}{处于社会主义初级阶段}{中国共产党的领导}{仍为世界最大发展中国家}

\SubQ 中国共产党领导是中国特色社会主义（\quad）。
\opts{最本质的特征}{重要特征之一}{基本特征}{一般特征}

\SubQ 道路问题是关系党的事业兴衰成败的（\quad）的问题。
\opts{第一位}{根本}{关键}{核心}

\SubQ "一个变与两个没有变"中，"变"的是（\quad）。
\opts{社会主要矛盾}{基本国情}{国际地位}{发展阶段}

\SubQ 我国社会主要矛盾已经转化为人民日益增长的（\quad）之间的矛盾。
\opts{物质文化需要和落后社会生产}{美好生活需要和不平衡不充分的发展}{物质文化需要和不够发达的生产}{美好生活需要和不够先进的生产力}

\SubQ "两个结合"指的是把马克思主义基本原理同中国具体实际相结合、同（\quad）相结合。
\opts{中华优秀传统文化}{中国革命文化}{社会主义先进文化}{世界优秀文明成果}

\SubQ 在"两个结合"中，马克思主义基本原理是（\quad）。
\opts{根脉}{魂脉}{血脉}{命脉}

\SubQ 在"两个结合"中，中华优秀传统文化是（\quad）。
\opts{魂脉}{根脉}{命脉}{血脉}

\SubQ "十个明确"在习近平新时代中国特色社会主义思想中的地位是（\quad）。
\opts{世界观和方法论}{基本方略}{最为核心关键的组成部分}{实践成果}

\SubQ "十四个坚持"是新时代坚持和发展中国特色社会主义的（\quad）。
\opts{基本方略}{世界观和方法论}{实践成果}{核心观点}

\SubQ "六个必须坚持"体现了习近平新时代中国特色社会主义思想的世界观和（\quad）。
\opts{方法论}{实践论}{认识论}{价值论}

\SubQ 习近平新时代中国特色社会主义思想实现了马克思主义中国化时代化（\quad）。
\opts{新的飞跃}{第二次飞跃}{第三次飞跃}{第四次飞跃}

\SubQ 中国共产党的自身优势中，居于第一位的能力是（\quad）。
\opts{思想引领力}{政治领导力}{群众组织力}{社会号召力}

\SubQ 民心是最大的政治，（\quad）。
\opts{决定事业兴衰成败}{决定改革成败}{决定发展方向}{决定历史进程}

\SubQ 共同富裕是中国特色社会主义的（\quad）。
\opts{根本要求}{本质要求}{基本要求}{重要要求}

\SubQ 全面深化改革的总目标是完善和发展中国特色社会主义制度、（\quad）。
\opts{实现国家治理现代化}{推进国家治理体系和治理能力现代化}{全面建成小康社会}{实现高质量发展}

\SubQ 新发展理念的内容是（\quad）。
\opts{创新、协调、绿色、开放、共享}{改革、发展、稳定、开放、共享}{创新、发展、绿色、开放、共享}{创新、协调、绿色、改革、共享}

\SubQ 在新发展理念中，注重解决发展动力问题的是（\quad）。
\opts{创新}{协调}{绿色}{开放}

\SubQ 在新发展理念中，注重解决人与自然和谐共生问题的是（\quad）。
\opts{创新}{协调}{绿色}{共享}

\SubQ 科技是第一生产力、人才是第一资源、（\quad）是第一动力。
\opts{改革}{开放}{创新}{发展}

\SubQ 人民民主是社会主义的（\quad）。
\opts{本质要求}{生命}{基本特征}{根本保证}

\SubQ 马克思主义认为，民主首先是、而且主要是一种（\quad）。
\opts{国家制度}{治理方式}{决策程序}{权利形式}

\SubQ 文化自信是更基础、更广泛、（\quad）的自信。
\opts{更深厚}{更持久}{更鲜明}{更强大}

\SubQ 意识形态关乎旗帜、关乎道路、关乎（\quad）。
\opts{国家政治安全}{经济发展}{社会稳定}{民族团结}

\SubQ 社会主义基本经济制度新概括中，作为所有制基础的是（\quad）。
\opts{"两个毫不动摇"}{按劳分配为主体}{社会主义市场经济体制}{公有制为主体}

\SubQ 新发展格局是（\quad）。
\opts{以国际大循环为主体}{以国内大循环为主体、国内国际双循环相互促进}{以国内国际双循环并重}{以国内单循环为主体}

\SubQ 生态文明建设应摆在（\quad）。
\opts{全局工作的突出位置}{经济工作的中心位置}{社会工作的首要位置}{政治工作的核心位置}

\vspace{10pt}
\examrule

% ============================================================
%  第二大题：多项选择题（10 题，每题 2 分，共 20 分）
% ============================================================
\BigQ{多项选择题（每题 2 分，共 20 分。每题有两个或两个以上正确选项，少选、多选、错选均不得分）}

\SubQ 习近平新时代中国特色社会主义思想创立的时代背景包括（\quad）。
\mopts{世界百年未有之大变局加速演进}{中华民族伟大复兴进入关键时期}{中国式现代化全面推进拓展}{科学社会主义在21世纪的中国焕发新的蓬勃生机}

\SubQ 推进中国式现代化必须牢牢把握的重大原则包括（\quad）。
\mopts{坚持和加强党的全面领导}{坚持中国特色社会主义道路}{坚持以人民为中心的发展思想}{坚持深化改革开放}

\SubQ 中国共产党的自身优势包括（\quad）。
\mopts{政治领导力}{思想引领力}{群众组织力}{社会号召力}

\SubQ 中国式现代化的中国特色包括（\quad）。
\mopts{人口规模巨大的现代化}{全体人民共同富裕的现代化}{物质文明和精神文明相协调的现代化}{人与自然和谐共生的现代化}

\SubQ 全过程人民民主的"四个相统一"包括（\quad）。
\mopts{过程民主和成果民主的统一}{程序民主和实质民主的统一}{直接民主和间接民主的统一}{人民民主和国家意志的统一}

\SubQ "五位一体"总体布局包括（\quad）。
\mopts{经济建设}{政治建设}{文化建设}{社会建设}

\SubQ 新发展理念的实践要求包括（\quad）。
\mopts{必须坚持创新在我国现代化建设全局中的核心地位}{必须正确处理局部和全局、当前和长远、重点和非重点的关系}{必须实现经济社会发展和生态环境保护协同共进}{必须坚持全民共享、全面共享、共建共享、渐进共享}

\SubQ 全面深化改革开放要坚持正确的方法论，主要包括（\quad）。
\mopts{增强系统性、整体性、协同性}{加强顶层设计和摸着石头过河相结合}{统筹改革发展稳定}{胆子要大，步子要稳}

\SubQ 中国特色社会主义文化的三个主要组成部分包括（\quad）。
\mopts{中华优秀传统文化}{革命文化}{社会主义先进文化}{外来文化}

\SubQ "一个变与两个没有变"中，"两个没有变"指的是（\quad）。
\mopts{我国仍处于并将长期处于社会主义初级阶段的基本国情没有变}{我国是世界最大发展中国家的国际地位没有变}{社会主要矛盾没有变}{发展阶段没有变}

\vspace{10pt}
\examrule

% ============================================================
%  第三大题：简答题（5 题，每题 10 分，共 50 分）
% ============================================================
\BigQ{简答题（每题 10 分，共 50 分。要求：在给出正确答案的基础上进行简要论述，否则不能得满分）}

\SubQ（10 分）试述习近平新时代中国特色社会主义思想创立的时代背景。

\vspace{8pt}
\SubQ（10 分）试述"两个确立"的决定性意义。

\vspace{8pt}
\SubQ（10 分）试述新时代的内涵（即中国特色社会主义新时代的五个"是"）。

\vspace{8pt}
\SubQ（10 分）试述推进中国式现代化必须牢牢把握的重大原则（五个方面），并说明这些原则与"四项基本原则"的关系。

\vspace{8pt}
\SubQ（10 分）试述新发展理念的科学内涵与实践要求。

\vspace{16pt}
\examrule
\vspace{4pt}
\begin{center}
  {\large\bfseries ——试题到此结束，请认真检查——}
\end{center}
\vspace{4pt}
\examrule

% ============================================================
%  参考答案
% ============================================================
\newpage
\setcounter{bigq}{0}

\begin{center}
  {\LARGE\bfseries 参考答案与评分标准}
\end{center}
\vspace{8pt}
\examrule

% -------- 第一大题参考答案 --------
\BigQ{单项选择题参考答案}

\begin{answerbox}
\renewcommand{\arraystretch}{1.4}
\begin{tabularx}{\linewidth}{@{}lX@{}}
1.\ A & 2.\ D \quad 3.\ A \quad 4.\ C \quad 5.\ A \quad 6.\ A \\
7.\ A & 8.\ B \quad 9.\ A \quad 10.\ B \quad 11.\ B \quad 12.\ C \\
13.\ A & 14.\ A \quad 15.\ A \quad 16.\ B \quad 17.\ A \quad 18.\ B \\
19.\ B & 20.\ A \quad 21.\ A \quad 22.\ C \quad 23.\ C \quad 24.\ B \\
25.\ A & 26.\ A \quad 27.\ A \quad 28.\ A \quad 29.\ B \quad 30.\ A \\
\end{tabularx}

\medskip
\noindent\textbf{评分说明}：每题 1 分，共 30 分。所选答案与参考答案完全一致方可得分，错选、多选、不选均不得分。
\end{answerbox}

% -------- 第二大题参考答案 --------
\BigQ{多项选择题参考答案}

\begin{answerbox}
\renewcommand{\arraystretch}{1.4}
\begin{tabularx}{\linewidth}{@{}lX@{}}
1.\ ABCD & 2.\ ABCD \quad 3.\ ABCD \quad 4.\ ABCD \quad 5.\ ABCD \\
6.\ ABCD & 7.\ ABCD \quad 8.\ ABCD \quad 9.\ ABC \quad 10.\ AB \\
\end{tabularx}

\medskip
\noindent\textbf{评分说明}：每题 2 分，共 20 分。每题有两个或两个以上正确选项，少选、多选、错选均不得分。
\end{answerbox}

% -------- 第三大题参考答案 --------
\BigQ{简答题参考答案}

\begin{answerbox}
\textbf{（1）试述习近平新时代中国特色社会主义思想创立的时代背景。（10 分）}

\textbf{【参考答案】}

习近平新时代中国特色社会主义思想创立的时代背景主要包括以下五个方面：

\textbf{第一，世界百年未有之大变局加速演进。} 当今世界正经历大发展大变革大调整时期，国际力量对比发生深刻变化，全球治理体系加速重塑，新一轮科技革命和产业变革深入发展，世界进入新的动荡变革期。这种深刻的国际格局变化，要求我们党对时代发展大势作出科学判断，回答好"世界怎么了、我们怎么办"的时代之问。

\textbf{第二，中华民族伟大复兴进入关键时期。} 党的十八大以来，中国特色社会主义进入新时代，全面建成小康社会取得伟大历史性成就，实现第一个百年奋斗目标，正向第二个百年奋斗目标迈进。中华民族伟大复兴战略全局与世界百年未有之大变局相互交织、相互激荡，对党的创新理论提出新的更高要求。

\textbf{第三，中国式现代化全面推进拓展。} 中国式现代化是在中国共产党领导下的社会主义现代化，既有各国现代化的共同特征，更有基于自己国情的中国特色。中国式现代化的全面推进拓展，为党的创新理论提供了坚实的实践基础。

\textbf{第四，科学社会主义在21世纪的中国焕发新的蓬勃生机。} 苏东剧变后，西方曾鼓吹"历史终结论"，但中国特色社会主义的成功实践，使科学社会主义在21世纪的中国焕发出新的蓬勃生机，彰显了社会主义制度的强大生命力。

\textbf{第五，中国共产党在革命性锻造中更加坚强有力。} 党的十八大以来，以习近平同志为核心的党中央以巨大的政治勇气和强烈的责任担当，推动党和国家事业取得历史性成就、发生历史性变革，党在革命性锻造中更加坚强有力，为这一思想的创立提供了坚强的政治保证和组织基础。

\textbf{【评分标准】} 五个方面每个方面 2 分，共 10 分。要点准确、表述规范即可得分；若只罗列要点不展开论述，每点扣 0.5 分。
\end{answerbox}

\begin{answerbox}
\textbf{（2）试述"两个确立"的决定性意义。（10 分）}

\textbf{【参考答案】}

"两个确立"是指：确立习近平同志党中央的核心、全党的核心地位，确立习近平新时代中国特色社会主义思想的指导地位。其决定性意义主要体现在以下几个方面：

\textbf{第一，"两个确立"是新时代取得一切成就的根本原因。} 党的十八大以来，党和国家事业取得历史性成就、发生历史性变革，最根本的原因在于有习近平总书记作为党中央的核心、全党的核心掌舵领航，在于有习近平新时代中国特色社会主义思想科学指引。这两个方面构成了新时代党和国家事业发展的根本政治保证和思想保证。

\textbf{第二，"两个确立"是战胜一切艰难险阻、应对一切不确定性的最大确定性。} 当今世界正经历百年未有之大变局，我国发展面临的风险挑战前所未有。有了坚强的领导核心和科学的理论指引，全党全国各族人民就有了"主心骨"，就能够凝聚起战胜一切困难、应对一切风险的强大力量。

\textbf{第三，"两个确立"是应对一切不确定性的最大底气。} 面对错综复杂的国内外形势和各种可以预见与难以预见的风险挑战，确立习近平同志党中央的核心、全党的核心地位，使全党有了团结统一的"定盘星"；确立习近平新时代中国特色社会主义思想的指导地位，使党和人民事业有了"指南针"。

\textbf{第四，"两个确立"是应对一切不确定性的最大保证。} 历史和现实都证明，有没有一个成熟的马克思主义执政党特别是领导核心，关系到党的兴衰成败，关系到社会主义事业的命运前途。"两个确立"反映了全党全军全国各族人民共同心愿，对新时代党和国家事业发展、对推进中华民族伟大复兴历史进程具有决定性意义。

\textbf{第五，要注意"两个确立"与"两个维护"的内在联系。} "两个确立"侧重思想认识层面的政治论断，"两个维护"侧重实践行动层面的政治要求，二者紧密联系、相互贯通。要紧密结合、一体领会，做到政治认同、思想认同、理论认同、情感认同。

\textbf{【评分标准】} 阐明"两个确立"内涵 2 分；从四个"最大"角度论述其决定性意义 6 分；阐述与"两个维护"的关系 2 分，共 10 分。要点准确、论述充分方可得满分。
\end{answerbox}

\begin{answerbox}
\textbf{（3）试述新时代的内涵（即中国特色社会主义新时代的五个"是"）。（10 分）}

\textbf{【参考答案】}

中国特色社会主义进入新时代，是我国发展新的历史方位。新时代的内涵主要体现在以下五个方面：

\textbf{第一，这个新时代，是承前启后、继往开来、在新的历史条件下继续夺取中国特色社会主义伟大胜利的时代。} 这一论断强调新时代的历史连续性和创新性，要求我们党在已有成就基础上，继续推进中国特色社会主义伟大事业，不断开创事业发展新局面。

\textbf{第二，是决胜全面建成小康社会、进而全面建设社会主义现代化强国的时代。} 这一论断明确了新时代的战略目标，即先决胜全面建成小康社会，实现第一个百年奋斗目标，再开启全面建设社会主义现代化国家新征程，向第二个百年奋斗目标进军。

\textbf{第三，是全国各族人民团结奋斗、不断创造美好生活、逐步实现全体人民共同富裕的时代。} 这一论断体现了新时代的人民性，强调以人民为中心的发展思想，把人民对美好生活的向往作为奋斗目标，扎实推进共同富裕。

\textbf{第四，是全体中华儿女勠力同心、奋力实现中华民族伟大复兴中国梦的时代。} 这一论断突出了新时代的使命担当，要求凝聚全体中华儿女的智慧和力量，为实现中华民族伟大复兴的中国梦而共同奋斗。

\textbf{第五，是我国不断为人类作出更大贡献的时代。} 这一论断彰显了新时代的天下情怀和国际担当，体现了中国作为负责任大国的胸怀与格局，推动构建人类命运共同体，为人类和平与发展事业作出新的更大贡献。

\textbf{【特别提示】} 注意区分"一个变与两个没有变"：变的是社会主要矛盾，没有变的是我国仍处于社会主义初级阶段的基本国情和仍是世界最大发展中国家的国际地位。新时代的伟大变革，在党史、新中国史、改革开放史、社会主义发展史、中华民族发展史上都具有里程碑意义。

\textbf{【评分标准】} 五个"是"每个方面 2 分，共 10 分。要点准确、表述规范即可得分；如能联系"一个变两个没有变"或"五大里程碑意义"加以补充，可酌情加分（但总分不超过 10 分）。
\end{answerbox}

\begin{answerbox}
\textbf{（4）试述推进中国式现代化必须牢牢把握的重大原则（五个方面），并说明这些原则与"四项基本原则"的关系。（10 分）}

\textbf{【参考答案】}

\textbf{一、五大重大原则的内容（5 分）}

推进中国式现代化必须牢牢把握以下五大重大原则：

\textbf{第一，坚持和加强党的全面领导。} 党的领导是中国特色社会主义最本质的特征，是中国式现代化的根本保证。只有坚持和加强党的全面领导，中国式现代化才能前景光明、繁荣兴盛。

\textbf{第二，坚持中国特色社会主义道路。} 道路决定命运，中国特色社会主义道路是党和人民在长期实践中开辟出来的正确道路，是实现社会主义现代化、创造人民美好生活的必由之路。

\textbf{第三，坚持以人民为中心的发展思想。} 中国式现代化是亿万人民自己的事业，必须坚持以人民为中心，发展为了人民、发展依靠人民、发展成果由人民共享，让现代化建设成果更多更公平惠及全体人民。

\textbf{第四，坚持深化改革开放。} 改革开放是决定当代中国命运的关键一招，也是决定中国式现代化成败的关键一招。要不断深化各领域改革，扩大高水平对外开放，为现代化建设注入强大动力。

\textbf{第五，坚持发扬斗争精神。} 推进中国式现代化是一项前无古人的开创性事业，必然会遇到各种可以预料和难以预料的风险挑战，必须增强全党全国各族人民的志气、骨气、底气，发扬斗争精神，增强斗争本领。

\textbf{二、五大原则与"四项基本原则"的关系（5 分）}

\textbf{第一，二者本质上是一脉相承、与时俱进的关系。} "四项基本原则"是立国之本，是党和国家生存发展的政治基石，即坚持社会主义道路、坚持人民民主专政、坚持中国共产党的领导、坚持马克思列宁主义毛泽东思想。五大重大原则是新时代推进中国式现代化必须遵循的根本原则，是在新的历史条件下对"四项基本原则"的继承、丰富和发展。

\textbf{第二，二者在核心要义上相互贯通。} 五大原则中，坚持和加强党的全面领导对应"四项基本原则"中的坚持中国共产党的领导；坚持中国特色社会主义道路对应坚持社会主义道路；坚持以人民为中心的发展思想体现人民民主专政的阶级本质和人民性；坚持深化改革开放和坚持发扬斗争精神，则是新时代对坚持马克思列宁主义毛泽东思想的实践展开。

\textbf{第三，二者在历史方位上相互衔接。} "四项基本原则"是邓小平同志在改革开放初期提出的，主要解决改革开放沿着正确方向前进的政治保证问题；五大重大原则是新时代新征程上推进中国式现代化的根本遵循，主要解决如何在新征程上行稳致远的问题，体现了党在不同历史阶段对坚持和发展中国特色社会主义规律认识的深化。

\textbf{【评分标准】} 五大原则每点 1 分，共 5 分；与"四项基本原则"关系从三个角度论述，每个角度 1\textasciitilde 2 分，共 5 分。要点准确、论述充分方可得满分。
\end{answerbox}

\begin{answerbox}
\textbf{（5）试述新发展理念的科学内涵与实践要求。（10 分）}

\textbf{【参考答案】}

新发展理念是习近平经济思想的重要内容，是新时代我国经济社会发展的指导原则。其科学内涵与实践要求如下：

\textbf{一、科学内涵（5 分）}

\textbf{第一，创新是引领发展的第一动力，创新发展注重的是解决发展动力问题。} 创新发展强调把创新摆在国家发展全局的核心位置，不断推进理论创新、制度创新、科技创新、文化创新等各方面创新，推动发展从要素驱动向创新驱动转变。

\textbf{第二，协调是持续健康发展的内在要求，协调发展注重的是解决发展不平衡问题。} 协调发展强调增强发展的整体性、协调性，正确处理局部和全局、当前和长远、重点和非重点的关系，促进区域协调发展、城乡协调发展、物质文明和精神文明协调发展。

\textbf{第三，绿色是永续发展的必要条件和人民对美好生活追求的重要体现，绿色发展注重的是解决人与自然和谐共生问题。} 绿色发展强调把生态文明建设摆在突出位置，加快形成节约资源和保护环境的空间格局、产业结构、生产方式、生活方式，实现经济社会发展和生态环境保护协同共进。

\textbf{第四，开放是国家繁荣发展的必由之路，开放发展注重的是解决发展内外联动问题。} 开放发展强调顺应我国经济深度融入世界经济的趋势，奉行互利共赢的开放战略，发展更高层次的开放型经济，构建广泛的利益共同体。

\textbf{第五，共享是中国特色社会主义的本质要求，共享发展注重的是解决社会公平正义问题。} 共享发展强调坚持全民共享、全面共享、共建共享、渐进共享，使全体人民在共建共享发展中有更多获得感，朝着共同富裕方向稳步前进。

\textbf{二、实践要求（5 分）}

\textbf{第一，必须坚持创新在我国现代化建设全局中的核心地位。} 让创新贯穿党和国家一切工作，全面提升创新能力和效率，把创新发展主动权牢牢掌握在自己手中，加快建设科技强国、制造强国、质量强国。

\textbf{第二，必须正确处理局部和全局、当前和长远、重点和非重点的关系。} 在发展中促进相对平衡，不断增强发展的整体性，统筹推进区域、城乡、物质文明与精神文明等各方面协调发展。

\textbf{第三，必须实现经济社会发展和生态环境保护协同共进。} 加快发展方式绿色转型，推动形成绿色低碳的生产方式和生活方式，建设人与自然和谐共生的现代化。

\textbf{第四，必须推动形成更大范围、更宽领域、更深层次对外开放格局。} 不断增强我国国际经济合作和竞争新优势，构建以国内大循环为主体、国内国际双循环相互促进的新发展格局。

\textbf{第五，必须坚持全民共享、全面共享、共建共享、渐进共享。} 不断推进全体人民共同富裕，正确处理效率和公平的关系，构建初次分配、再分配、第三次分配协调配套的基础性制度安排。

\textbf{【评分标准】} 科学内涵五个方面每点 1 分，共 5 分；实践要求五个方面每点 1 分，共 5 分。要点准确、表述规范即可得分；如能联系"完整、准确、全面贯彻新发展理念"加以补充，可酌情加分（但总分不超过 10 分）。
\end{answerbox}

\vspace{12pt}
\examrule
\begin{center}
  \textit{参考答案仅供参考，评分时应酌情给分，步骤正确可得过程分。}
\end{center}

\end{document}
